%%%%%%%%%%%%%%%%%%%%%%%%%%%%%%%%%%%%%%%%%%%%%%%%%%%%%%%%%%%%%%%%%%%%%%%%%%%%%
%                                                                           %
%  Helper file to generate a single chapter in the tuProlog documentation.  %
%                                                                           %
%   Andrea Omicini, 2/8/2010                                                %
%                                                                           %
%%%%%%%%%%%%%%%%%%%%%%%%%%%%%%%%%%%%%%%%%%%%%%%%%%%%%%%%%%%%%%%%%%%%%%%%%%%%%

\documentclass[11pt]{report}
%%%%%%%%%%%%%%%%%%%%%%%%%%%%%%%%%%%%%%%%%%%%%%%%%%%%%%%%%%%%%%%%%%%%%%%%%
\usepackage[pdftex]{graphicx}
\DeclareGraphicsExtensions{.png, .jpg, .eps, .pdf, .tif}
%%%%%%%%%%%%%%%%%%%%%%%%%%%%%%%%%%%%%%%%%%%%%%%%%%%%%%%%%%%%%%%%%%%%%%%%%
\newcommand\xa[1]{\appendixname~\ref{app:#1}}
\newcommand\labelsec[1]{\label{sec:#1}}
\newcommand\xs[1]{\sectionname~\ref{sec:#1}}
\newcommand\xsp[1]{\sectionname~\ref{sec:#1} \onpagename~\pageref{sec:#1}}
\newcommand\labelssec[1]{\label{ssec:#1}}
\newcommand\xss[1]{\subsectionname~\ref{ssec:#1}}
\newcommand\xssp[1]{\subsectionname~\ref{ssec:#1} \onpagename~\pageref{ssec:#1}}
\newcommand\labelsssec[1]{\label{sssec:#1}}
\newcommand\xsss[1]{\subsectionname~\ref{sssec:#1}}
\newcommand\xsssp[1]{\subsectionname~\ref{sssec:#1} \onpagename~\pageref{sssec:#1}}
\newcommand\labelfig[1]{\label{fig:#1}}
\newcommand\xf[1]{\figurename~\ref{fig:#1}}
\newcommand\xff[2]{\figurenames~\ref{fig:#1}~and~\ref{fig:#2}}
\newcommand\xfp[1]{\figurename~\ref{fig:#1} \onpagename~\pageref{fig:#1}}
\newcommand\labeltab[1]{\label{tb:#1}}
\newcommand\xt[1]{\tablename~\ref{tb:#1}}
\newcommand\xtt[2]{\tablenames~\ref{tb:#1}~and~\ref{ab:#2}}
\newcommand\xtp[1]{\tablename~\ref{tb:#1} \onpagename~\pageref{tb:#1}}
\newcommand\labelenum[1]{\label{enum:#1}}
\newcommand\xen[1]{(\ref{enum:#1})}
\newcommand\xenp[1]{(\ref{enum:#1}) \onpagename~\pageref{enum:#1}}
%******************************************************************************%
\newcommand\bti[1]{\texttt{\textbf{#1}}}
\newcommand\bt[1]{\texttt{#1}}
\newcommand\template[1]{\textit{Template: }\texttt{#1}}
\newcommand\ttit[1]{\texttt{\textit{#1}}}
%******************************************************************************%
\newcommand{\LIA}{\mbox{\textsf{LIA}}}
\newcommand{\aclt}{\mbox{$\mathcal{ACLT}$}}
\newcommand{\respect}{\mbox{\sf{{R}e{S}pec{T}}}}
\newcommand{\luce}{\mbox{\sf{{L}u{C}e}}}
\newcommand{\tucson}{\mbox{\sf{{T}u{CS}o{N}}}}
\newcommand{\alice}{\mbox{\sf{{aliCE}}}}
\newcommand{\tuprolog}{\mbox{{\sf{tu}}Prolog}}
\newcommand{\apice}{\mbox{\sf{{APICe}}}}
%******************************************************************************%
\newcommand\version[1]{\mbox{Document revision: #1}}
\newcommand\approvedby[1]{\mbox{Approved by: #1}}
\newcommand\receivedby[1]{\mbox{Received by: #1}}
\newcommand\creationdate[1]{\mbox{Creation date: #1}}
\newcommand\lastchangesdate[1]{\mbox{Last Changes date: #1}}
\newcommand\noa[2]{\noindent\emph{Note of the author (#1): }#2\\\\}
\newcommand\logo{
    \begin{figure}[tp]
        \begin{center}
            \includegraphics[width=2cm]{common/logo}
        \end{center}
\end{figure}
}
%******************************************************************************%
\newcommand{\classname}[1]{\texttt{#1}}
\newcommand{\varname}[1]{\texttt{#1}}
\newcommand{\predicate}[1]{\texttt{#1}}
\newcommand{\code}[1]{\texttt{#1}}
\newcommand{\keycap}[1]{\textbf{#1}}
\newcommand{\guibutton}[1]{\textsf{#1}}
\newcommand{\userinput}[1]{\texttt{#1}}
\newcommand{\figref}[1]{\figurename~\ref{#1}}
%******************************************************************************%
%\title{{\huge{\bf{\tuprolog{} Guide\\\mbox{ }\\}}}
%        \tuprolog{} version: 1.2.1\\\mbox{ }\\
%        \tuprolog{} IDE version: 1.1.0\\\mbox{ }\\
%{\small{
%    \version{001}\\
%    \creationdate{2002-09-17}\\
%    \lastchangesdate{2004-09-27}\\
%    % \receivedby{aricci}\\
%    % \approvedby{aricci}\\
%    }}
%}
%
%\author{ \mbox{ }\\DEIS, Universit\`{a} di Bologna a Cesena, Italy
%}

%\date{}

\usepackage{supertabular}
\usepackage[bookmarks=true,bookmarksopen=true,bookmarksnumbered]{hyperref}

\begin{document}

%\logo

%\maketitle

%\tableofcontents
\setcounter{chapter}{0}
%=======================================================================
\chapter{What is \tuprolog{}}
\label{what-is}
%=======================================================================

\tuprolog{} is an open-source, light-weight Prolog framework for distributed applications and infrastructures, released under the LGPL license, available from \url{http://tuprolog.apice.unibo.it}.

Originally developed in/upon Java, which still remains the main reference platform,
\tuprolog{} is currently available for several platforms/environments:
\begin{itemize}
  \item plain JavaSE;
  \item Eclipse plugin;
  \item Android;
  \item Microsoft .NET.
\end{itemize}
%
While they all share the same core and libraries, the latter features an \emph{ad hoc} library which extends the multi-paradigm approach to virtually any language available on the .NET platform (more on this in Section \xch{dotnet}).

Unlike most Prolog programming environments, aimed at providing a very efficient (yet monolithic) stand-alone Prolog system, \tuprolog{} is explicitly designed to be \emph{minimal}, dynamically \emph{configurable}, straightforwardly \emph{integrated} with Java and .NET so as to naturally support multi-paradigm/multi-language programming (MPP), and \emph{easily deployable}.

\textit{Minimality} means that its core contains only the Prolog engine essentials -- roughly speaking, the resolution engine and some related basic mechanisms -- for as little as 155KB: any other feature is implemented in \textit{libraries}.
%
So, each user can customize his/her prolog system to fit his/her own needs, and no more:
this is what we mean by \tuprolog{} \textit{configurability}---the necessary counterpart of minimality.

Libraries provide packages of predicates, functors and operators, and can be loaded and unloaded in a \tuprolog{} engine both statically and dynamically.
%
Several standard libraries are included in the \tuprolog{} distribution, and are loaded by default in the standard \tuprolog{} configuration; however, users can easily develop their own libraries either in several ways -- just pure Prolog, just pure Java\footnote{The .NET version of \tuprolog{} supports other languages available on the .NET platform: more on this topic in Section \xch{dotnet}}, or a mix of the two --, as we will discuss in \xch{howto-develop-libraries}.

\textit{Multi-paradigm programming} is another key feature of \tuprolog{}.
%
In fact, the \tuprolog{} design was intentionally calibrated from the early stages to support a straightforward, pervasive, multi-language/multi-paradigm integration, so as to enable users to:
\begin{itemize}
  \item using any Java\footnote{For the .NET version: any .NET class, library, object, etc.} class, library, object \emph{directly from the Prolog code}
  (\xch{java-library}) with no need of pre-declarations, awkward syntax, etc., with full support of parameter passing from the two worlds, yet leaving the two languages and computational models totally separate so as to preserve \emph{a priori} their own semantics---thus bringing the power of the object-oriented platform (e.g. Java Swing, JDBC, etc) to the Prolog world for free;

  \item using any Prolog engine \emph{directly from the Java/.NET code} as one would
   do with any other Java libraries/.NET assemblies, again with full support of parameter passing from the two worlds in a non-intrusive, simple way that does not alter any semantics---thus bringing the power of logic programming into virtually \emph{any} Java/.NET application;

  \item augmenting Prolog by defining new libraries \xch{howto-develop-libraries})
  either in Prolog, or in the object-oriented language of the selected platform (again, with a straightforward, easy-to-use approach based on reflection which avoids any pre-declaration, language-to-language mapping, etc), or in a mix of both;

  \item augmenting Java\footnote{This feature is currently available only in the Java version: a suitable extension to the .NET platform is under study.} by defining new Java methods in Prolog (the so-called `P@J' framework---\xch{PJ}), which exploits reflection and type inference to provide the user with an easy-to-use way to implement Java methods declaratively.
\end{itemize}

Last but not least, \textit{easy deployability} means that the installation requirements are minimal, and that the installation procedure is in most cases\footnote{Exceptions are the Eclipse plugin and the Android versions, which need to be installed as required by the hosting platforms.} as simple as copying one archive to the desired folder.
%
Coherently, a Java-based installation requires only a suitable Java Virtual Machine, and `installing' is just copying a single JAR file somewhere---for as much as 474KB of disk usage (yes, minimality is not just a claim here).
%
Of course, other components can be added (documentation, extra libraries, sources..), but are not necessary for a standard everyday use.
%
The file size is quite similar for the Android platform -- the single APK archive is 234KB -- although an Android-compliant install is performed due to Android requirements.
%
The install process is also quite the same on the .NET platform, although the files are slightly larger.
%
The Eclipse platform also requires a different procedure, since plugin installation have to conform to the requirements of the Eclipse plugin manager: consequently, an update site was set up, where the \tuprolog{} plugin is available as an Eclipse feature. Due to these constraints, file size increases to 1.5MB.


In order to manage all these platforms in a uniform way, a suitable \emph{version numbering scheme} was recently introduced:
\begin{itemize}
  \item the fist two digits represent the engine version;
  \item the last (third) digit is platform-specific and accounts for version differences which do not impact on the Prolog engine -- that is, on the
      \tuprolog{} behaviour -- but simply on graphical aspects or platform-specific
      issues or bugs.
\end{itemize}
%
So, as long as the first two digits are the same, a \tuprolog{} application is guaranteed to behave identically on any supported platform.


Finally, \tuprolog{} also supports \textit{interoperability} with both Internet standard patterns (such as TCP/IP, RMI, CORBA) and coordination models and languages.
%
The latter aspect, in particular, is currently developed in the context of the \tucson{} coordination infrastructure \cite{tucson-aamas99,respect-scico2001}, which provides logic-based, programmable tuple spaces (called \emph{tuple centres}) as the coordination media for distributed processes and agents.\footnote{An alternative infrastructure, \luce{} \cite{luce-aamas2001}, developed the same approach in a location-unaware fashion: this infrastructure is currently no longer supported.}



%\bibliography{common/references}
%\bibliographystyle{plain}

\end{document}
